\documentclass[11pt]{article}
\usepackage[margin=1in]{geometry}
\usepackage{hyperref}
\usepackage{longtable}
\usepackage{array}
\usepackage{enumitem}
\usepackage{graphicx}
\usepackage{listings}

\title{Code Origin Detector Project Report}
\author{Project Overview}
\date{\today}

\begin{document}

\maketitle
\tableofcontents
\newpage

\section{Overview}
\label{sec:overview}
The Code Origin Detector is a research-focused toolkit that estimates whether code was written by human developers or generated by AI systems. It combines language-aware heuristics, handcrafted features, and probabilistic models to deliver verdicts with supporting explanations.\footnote{See \texttt{README.md} for the project summary.}

\section{Repository Structure}
\label{sec:structure}
The project repository is organized to separate backend detection logic, frontend visualizations, documentation, scripts, and experiments:
\begin{itemize}[leftmargin=*]
  \item Root metadata includes \texttt{pyproject.toml} for packaging and \texttt{requirements.txt} for CLI dependencies.
  \item \texttt{docs/} stores design notes, prompt catalogs, reports, and example walkthroughs.
  \item \texttt{frontend/} contains the React~19 and Vite~7 dashboard implementation with Tailwind CSS styling.
  \item \texttt{notebooks/} hosts exploratory data analysis and modeling notebooks.
  \item \texttt{scripts/} provides automation for data collection and preprocessing.
  \item \texttt{src/detector/} implements the Python package with CLI entrypoints, heuristics, feature extractors, models, and utilities.
  \item \texttt{tests/} contains Python unit and integration suites for the backend.
\end{itemize}

\section{Backend Architecture}
\label{sec:backend}
The backend exposes a Typer-based CLI with commands for prediction, explanations, and evaluation. The \texttt{predict} command loads a configurable model, applies it across files or directories, and renders tabular or JSON outputs. The \texttt{explain} command focuses on a single file and lists top rationale signals, while the \texttt{eval} command benchmarks a manifest of labeled files and prints aggregate metrics and confusion matrices.\footnote{See \texttt{src/detector/cli.py}.}

Central to inference is the \texttt{Predictor} class. It constructs a \texttt{FeaturePipeline} that chains stylometric, Python AST, and JavaScript heuristics extractors, enforces a shared \texttt{FeatureIndex}, and loads either the rule-based heuristic model or serialized scikit-learn estimators. Predictions include detected language, probability of AI authorship, confidence, triggered heuristics, and merged explanations. Batch predictions walk the filesystem via inclusion/exclusion globs, while evaluation routines compute accuracy, balanced accuracy, macro-F1, AUROC, and confusion matrices using scikit-learn utilities.\footnote{See \texttt{src/detector/inference.py}.}

\section{Feature Extraction}
\label{sec:features}
Feature extraction is modular. The \texttt{FeaturePipeline} iterates over registered extractors that declare language applicability, namespaces feature keys, and returns a combined dictionary for downstream models. The stylometry extractor yields line-level statistics such as counts, comment ratios, indentation, and whitespace density, supporting cross-language stylistic analysis.\footnote{See \texttt{src/detector/featurizers/base.py} and \texttt{src/detector/featurizers/stylometry.py}.}

Python- and JavaScript-specific extractors extend the pipeline (not shown) with AST-driven metrics including identifier usage, comprehension counts, function lengths, and idiom prevalence, enabling richer heuristics and model features.

\section{Heuristic Signals}
\label{sec:heuristics}
Language-specific heuristic rules interpret feature dictionaries to produce \texttt{HeuristicSignal} records with identifiers, scores, human-readable messages, and evidence snippets. For Python, rules flag unusually generic identifiers, complex functions lacking comments, reliance on loops over comprehensions, and broad exception handlers without logging. Each rule emits weighted scores that inform the heuristic baseline and explanation surfaces.\footnote{See \texttt{src/detector/heuristics/\allowbreak python\_rules.py}.}

\section{Modeling and Evaluation}
\label{sec:modeling}
The \texttt{HeuristicModel} acts as the default detector, combining rule outputs into probability estimates. When provided, serialized scikit-learn models (e.g., logistic regression, random forests) are loaded alongside the saved feature index to ensure feature alignment. Evaluation on labeled manifests produces metrics (macro-F1, balanced accuracy, accuracy, AUROC) and confusion matrices for quality tracking.\footnote{See \texttt{src/detector/inference.py}.}

\section{Frontend Experience}
\label{sec:frontend}
The React dashboard mirrors CLI functionality with added usability features: drag-and-drop uploads, curated examples, configurable model profiles, heuristic overlays, verdict badges, confidence meters, and persistence of prediction settings. Tailwind design tokens unify palettes, spacing, and motion-safe focus states. The development workflow uses npm scripts for linting, testing, building, and formatting.\footnote{See \texttt{README.md}.}

\section{Testing and Quality}
\label{sec:quality}
Backend quality gates rely on Ruff, mypy, and pytest, while the frontend leverages ESLint, Jest with Testing Library, Vite build checks, and Prettier formatting. These checks ensure consistency across Python and TypeScript components.\footnote{See \texttt{README.md}.}

\section{Data Roadmap}
\label{sec:data}
Planned datasets include curated human-authored corpora sampled from vetted open-source commits and AI-generated corpora produced via scripted prompts. Deduplication leverages SHA-256 hashing and token shingles, with repository-level splits and temporal hold-outs. Metadata schemas in \texttt{data/metadata/schema.json} support reproducible ingestion.\footnote{See \texttt{README.md}.}

\section{Responsible Use}
\label{sec:responsible}
The detector provides advisory signals rather than definitive judgments. False positives and negatives remain possible, especially for boilerplate or expert-like generated code. Outputs are intended to guide manual review rather than serve as automated policy gates.\footnote{See \texttt{README.md}.}

\section{Mermaid Diagram Codes}
\label{sec:mermaid}
The following Mermaid definitions capture high-level system views for visualization tooling.

\subsection{Detection Pipeline}
\begin{verbatim}
```mermaid
flowchart TD
    A[Source Files] --> B[Language Detection]
    B --> C[Feature Pipeline]
    C --> D[Heuristic Rules]
    C --> E[Statistical Model]
    D --> F[Probability Aggregation]
    E --> F
    F --> G[Predictions + Explanations]
```
\end{verbatim}

\subsection{Repository Responsibilities}
\begin{verbatim}
```mermaid
mindmap
  root((Code Origin Detector))
    Backend
      CLI
      Feature Pipeline
      Heuristics
      Models
    Frontend
      React Dashboard
      Heuristic Overlays
      Result History
    Data
      Collection Scripts
      Deduplication
      Manifests
```
\end{verbatim}

\end{document}
